\documentclass[11pt,letterpaper]{article}
% A diagnostic for the HTML build: an ordinary article using only standard
% packages, touching the constructs a tex4ht + MathJax conversion tends to get
% wrong. Each (sub)section is one group of TESTING.md's checklist; compare
% test.pdf with the HTML section by section.

\usepackage[T1]{fontenc}
\usepackage{lmodern}
\usepackage{amsmath, amssymb, amsthm, mathtools, bm, mathrsfs}
\usepackage{booktabs}
\usepackage{tikz}
\usepackage[colorlinks]{hyperref}
\usepackage[nameinlink]{cleveref}

\newtheorem{theorem}{Theorem}[section]
\newtheorem{lemma}[theorem]{Lemma}
\newtheorem{corollary}[theorem]{Corollary}
\theoremstyle{definition}
\newtheorem{definition}[theorem]{Definition}
\newtheorem{example}[theorem]{Example}
\theoremstyle{remark}
\newtheorem*{remark}{Remark}

\DeclareMathOperator{\Var}{Var}
\DeclareMathOperator{\tr}{tr}
\DeclareMathOperator*{\argmax}{arg\,max}
\DeclarePairedDelimiter{\abs}{\lvert}{\rvert}
\DeclarePairedDelimiter{\norm}{\lVert}{\rVert}
\newcommand{\E}{\mathbb{E}}
\newcommand{\R}{\mathbb{R}}
\newcommand{\mr}{\mathrm}
\newcommand{\KL}{D_\mathrm{KL}}

\title{Lorem ipsum dolor sit amet}
\author{consectetur adipiscing elit}
\date{sed do eiusmod}

\begin{document}
\maketitle
\begin{abstract}
    Lorem ipsum dolor sit amet, consectetur adipiscing elit, sed do eiusmod tempor incididunt ut labore et dolore magna aliqua.
    Ut enim ad minim veniam, quis nostrud exercitation ullamco laboris nisi ut aliquip ex ea commodo consequat.
\end{abstract}
\tableofcontents


\section{Text}

\subsection{Type}

Lorem ipsum dolor sit amet, consectetur adipiscing elit, sed do eiusmod tempor incididunt ut labore et dolore magna aliqua.
Ligatures: fi, fl, ff, ffi, ffl; dashes: 1--2, a---b; quotes: ``double'' and `single'; accents: na\"ive, caf\'e, \v{C}ech; ellipsis\ldots{} and a tie in Fig.~\ref{fig:tikz}.
Styles: \emph{emphasis with \emph{nested} emphasis}, \textbf{bold}, \textit{\textbf{bold italic}}, \textsc{Small Caps}, \texttt{typewriter}, \textsf{sans serif}, \underline{underline}, and a footnote.\footnote{A footnote with math, $\sum_{i=1}^n x_i$, and a \href{https://example.com}{link}.}
A URL: \url{https://example.com/some_path?q=1&r=2}; inline verbatim: \verb|\frac{a}{b} % not a comment|.

Ut enim ad minim veniam, quis nostrud exercitation ullamco laboris nisi ut aliquip ex ea commodo consequat.
This paragraph is indented in the PDF.

\noindent This one is not.
\hfill Right-aligned text.

\subsection{Lists}

\begin{enumerate}
    \item First item, with inline math $a^2 + b^2 = c^2$.
    \item Second item, with a nested list:
    \begin{enumerate}
        \item nested (a);
        \item nested (b), with a display
        \[ x = \frac{-b \pm \sqrt{b^2 - 4ac}}{2a}. \]
    \end{enumerate}
    \item[$\star$] A custom label.
\end{enumerate}
\begin{itemize}
    \item Bullet.
    \item Bullet with a second paragraph.

    The second paragraph.
\end{itemize}
\begin{description}
    \item[Term] Its description.
    \item[Another term] Another description.
\end{description}

\subsection{Quotations and verbatim}

\begin{quote}
    A quotation. Duis aute irure dolor in reprehenderit in voluptate velit esse cillum dolore eu fugiat nulla pariatur.
\end{quote}
\begin{verbatim}
def f(x):
    return x ** 2   # verbatim block
\end{verbatim}


\section{Inline math}

\subsection{Alphabets}

Italic $x, y, \alpha, \beta, \Gamma, \varGamma, \ell, \imath, \hbar, \aleph$;
upright $\mathrm{d}x$, $\mathrm{e}^{i\pi}$, $\Gamma(z)$;
bold $\mathbf{x}, \mathbf{A}, \bm{x}, \bm{\beta}, \bm{\Sigma}, \bm{\ell}, \boldsymbol{\mu}$;
blackboard $\mathbb{R}^n, \E[X], \mathbb{1}$;
calligraphic $\mathcal{L}, \mathcal{F}$ vs script $\mathscr{L}, \mathscr{F}$;
fraktur $\mathfrak{g}$; sans $\mathsf{T}$; typewriter $\mathtt{x}$; italic $\mathit{diff}$;
text in math $\text{if } x > 0 \text{ and ff/fi}$ and $\operatorname{diff}(x)$, $\mathrm{ff}$.

\subsection{Operators}

$\sin x$, $\sin(x)$, $\sin^2 x$, $\log n$, $\log(n)$, $\exp\{x\}$, $\det A$, $\tr(A)$, $\Var(X)$, $\Var(X)^2$, $\E[X]^2$, $\Pr(A \mid B)$,
$\max_{i} a_i$, $\min\{a, b\}$, $\sup_{x \in A} f(x)$, $\lim_{n\to\infty} a_n$, $\limsup_{n} a_n$, $\liminf a_n$, $\argmax_\theta \ell(\theta)$,
$\sum_{k=1}^n k$, $\prod_i p_i$, $\int_0^1 f(x)\,dx$, $\oint_C$, $\iint_D$, $\bigcup_i A_i$, $\bigcap_{i} A_i$,
$a \bmod b$, $a \equiv b \pmod{n}$, $\gcd(a, b)$.

\subsection{Relations, binary operators and spacing}

$f\colon A \to B$, $x \mapsto x^2$, $\{x : x > 0\}$, $\{x \mid x > 0\}$, $\{x \in \R \;\big|\; x > 0\}$,
$a\,b\:c\;d\!e\quad f\qquad g$,
$-1$, $+x$, $(-x)$, $a - b$, $a \cdot b$, $a \times b$, $a \circ b$, $a \cup b$, $\neg p$, $p \wedge q$, $\forall x \exists y$,
$1, 2, \dots, n$, $a_1 + \dots + a_n$, $a_1 + \cdots + a_n$, $x \in \varnothing$, $x \notin \emptyset$,
$x \neq y$, $x \not= y$, $x \not\in A$, $a \leq b$, $a \leqslant b$, $a \nleq b$, $A \subseteq B$, $A \subsetneq B$, $P \preceq Q$, $P \succcurlyeq Q$, $f \propto g$, $X \sim N(0, 1)$, $x \approx y$, $a \coloneqq b$, $u \parallel v$, $\|v\|$, $X \perp\!\!\!\perp Y$.

\subsection{Scripts and primes}

$x^2$, $x_i^2$, ${x_i}^2$, $x^{y^z}$, $x_{i_j}$, $e^{-x^2/2}$, $e^{-\frac{x^2}{2}}$, ${}^n C_k$, $f'(x)$, $f''(x)$, $f^{(n)}$, $x'_i$, $(x')^2$, $A^\top$, $A^{-1}$, $A^\dagger$, $x^*$, $\bar{x}^2$,
$x_\mathrm{op}$, $x_\mr{op}$, $\KL(p \parallel q)$.

\subsection{Accents}

$\hat{\theta}$, $\hat{\bm\beta}$, $\widehat{\theta}$, $\widehat{xyz}$, $\tilde{x}$, $\widetilde{xyz}$, $\bar{x}$, $\overline{x + y}$, $\underline{x}$, $\vec{v}$, $\overrightarrow{AB}$, $\dot{x}$, $\ddot{x}$, $\acute{a}$, $\breve{a}$, $\check{a}$, $\mathring{A}$, $\hat{\imath}$.

\subsection{Fractions, binomials and roots}

$\frac{a}{b}$, $\tfrac{1}{2}$, $\dfrac{1}{2}$, $\frac{1}{1 + \frac{1}{x}}$, $\binom{n}{k}$, $\tbinom{n}{k}$, $\sqrt{x}$, $\sqrt[3]{x}$, $\sqrt{\frac{a}{b}}$, $\left(\begin{smallmatrix} a & b \\ c & d \end{smallmatrix}\right)$.

\subsection{Delimiters}

$(\frac{a}{b})$, $\left(\frac{a}{b}\right)$, $\bigl(x\bigr)$, $\Bigl[x\Bigr]$, $\biggl\{x\biggr\}$, $\Biggl\langle x\Biggr\rangle$, $\Bigl(\bigl(x\bigr)\Bigr)$, $\abs{x}$, $\abs*{\frac{x}{2}}$, $\abs[\big]{x}$, $\norm{x}_2$, $\lvert x \rvert$, $|x|$, $\lfloor x \rfloor$, $\lceil x \rceil$, $\langle u, v \rangle$, $\left. \frac{df}{dx} \right|_{x=0}$, $\left\{ x \;\middle|\; x > 0 \right\}$, $P(A \bigm| B)$.

\subsection{Math styles}

$\textstyle\sum_{i=1}^n$ vs $\displaystyle\sum_{i=1}^n$, and $x^{\scriptstyle a} x^{\scriptscriptstyle a}$.


\section{Display math}

\subsection{Large operators and limits}

\[ \sum_{k=1}^{\infty} \frac{1}{k^2} = \frac{\pi^2}{6}, \qquad \prod_{p \text{ prime}} \frac{1}{1 - p^{-s}} = \zeta(s), \qquad \int_{-\infty}^{\infty} e^{-x^2}\,dx = \sqrt{\pi}. \]
\[ \lim_{n \to \infty} \Bigl(1 + \frac{x}{n}\Bigr)^n = e^x, \quad \limsup_{n \to \infty} a_n, \quad \max_{\substack{1 \leq i \leq n \\ i \text{ odd}}} a_i, \quad \argmax_{\theta \in \Theta} \ell(\theta), \quad \textstyle\sum_{i=1}^n a_i, \quad \int\limits_0^1, \quad \sum\nolimits_k. \]

\subsection{Braces, arrows and boxes}

\[ \underbrace{a + \cdots + a}_{n \text{ times}} = \overbrace{n \cdot a}^{\text{product}}, \quad A \overset{\text{def}}{=} B, \quad X_n \xrightarrow[n \to \infty]{\text{d}} X, \quad x \xleftarrow{f} y, \quad \boxed{E = mc^2}. \]

\subsection{Big delimiters}

\[ \left( \sum_{i=1}^n \frac{x_i}{\sqrt{1 + x_i^2}} \right)^2 \leq \left\lVert \frac{\bm x}{\norm{\bm x}} \right\rVert \cdot \Biggl[ \int_0^1 \biggl( f(t) - \frac{1}{2} \biggr)^2 dt \Biggr], \qquad \left. \frac{\partial^2 f}{\partial x\, \partial y} \right|_{(0,0)}. \]

\subsection{Fractions and text}

\[ x = a_0 + \cfrac{1}{a_1 + \cfrac{1}{a_2 + \cfrac{1}{a_3 + \dotsb}}}, \qquad \genfrac{[}{]}{0pt}{}{n}{k}. \]
\[ f(x) = x^2 \quad \text{for all } x \in \R, \text{ with $x$ in text}, \qquad \phantom{xx} y, \qquad a \mathbin{\#} b \mathrel{\#} c. \]

\subsection{Numbered equations}

Referenced as \eqref{eq:numbered}, \cref{eq:split} and Equation~\ref{eq:split}:
\begin{equation}\label{eq:numbered}
    \Var(X) = \E\bigl[(X - \E[X])^2\bigr] = \E[X^2] - \E[X]^2.
\end{equation}
\begin{equation}\label{eq:split}
\begin{split}
    (a + b)^2 &= (a + b)(a + b) \\
              &= a^2 + 2ab + b^2.
\end{split}
\end{equation}

\subsection{Cases}

\[ \abs{x} = \begin{cases} x & \text{if } x \geq 0, \\ -x & \text{otherwise}, \end{cases} \qquad f(x) = \begin{dcases} \frac{1}{x} & x > 0 \\ 0 & x \leq 0. \end{dcases} \]

\subsection{Alignments}

An align with a note in the middle, a tag and a gap:
\begin{align}
    \norm{x + y}^2 &= \norm{x}^2 + 2\langle x, y \rangle + \norm{y}^2 \notag \\
                   &\leq \norm{x}^2 + 2\norm{x}\norm{y} + \norm{y}^2 \label{eq:cs} \\
    \intertext{by Cauchy--Schwarz, and so}
                   &= (\norm{x} + \norm{y})^2 \tag{$\ast$} \\[6pt]
    a &= b & c &= d
\end{align}
An align* whose lines start with a relation:
\begin{align*}
    \E[X] &= \int x\,dP \\
          &= \sum_{k} k \, P(X = k) \\
          &\quad + \text{a long continuation term}.
\end{align*}
Gather*, multline and alignat:
\begin{gather*}
    a = b + c \\
    \frac{a}{b} = \frac{c}{d}
\end{gather*}
\begin{multline}
    p(x) = a_0 + a_1 x + a_2 x^2 + a_3 x^3 + a_4 x^4 + a_5 x^5 + a_6 x^6 \\
    + a_7 x^7 + a_8 x^8 + a_9 x^9 + a_{10} x^{10}.
\end{multline}
\begin{alignat}{2}
    x &= y + z \qquad && \text{by definition} \\
    y &= 2z         && \text{by assumption}
\end{alignat}

\subsection{Matrices and arrays}

\[ \begin{pmatrix} a_{11} & \cdots & a_{1n} \\ \vdots & \ddots & \vdots \\ a_{m1} & \cdots & a_{mn} \end{pmatrix}, \quad \begin{bmatrix} 1 & 0 \\ 0 & 1 \end{bmatrix}, \quad \det \begin{vmatrix} a & b \\ c & d \end{vmatrix}, \quad \left[ \begin{array}{cc|c} 1 & 2 & 3 \\ \hline 4 & 5 & 6 \end{array} \right], \quad \begin{pmatrix*}[r] -1 & 2 \\ 3 & -4 \end{pmatrix*}. \]

\subsection{Wide displays}

\[ \Pr\Bigl( \Bigl\lvert \frac{1}{n} \sum_{i=1}^n X_i - \mu \Bigr\rvert > \varepsilon \Bigr) \leq 2 \exp\Bigl( -\frac{2 n^2 \varepsilon^2}{\sum_{i=1}^n (b_i - a_i)^2} \Bigr) \quad \text{for all } \varepsilon > 0 \text{ and } n \geq 1, \text{ by Hoeffding's inequality}. \]


\section{Theorems and \texorpdfstring{$L^2$}{L2} in a heading}

\begin{definition}[Inner product]
    An \emph{inner product} on a real vector space $V$ is a map $\langle \cdot, \cdot \rangle \colon V \times V \to \R$ that is bilinear, symmetric and positive definite.
\end{definition}
\begin{theorem}[Cauchy--Schwarz]\label{thm:cs}
    For all $x, y \in V$, $\abs{\langle x, y \rangle} \leq \norm{x}\norm{y}$.
\end{theorem}
\begin{proof}
    If $y = 0$ there is nothing to show. Otherwise, for all $t \in \R$,
    \[ 0 \leq \norm{x - ty}^2 = \norm{x}^2 - 2t\langle x, y \rangle + t^2 \norm{y}^2, \]
    and taking $t = \langle x, y \rangle / \norm{y}^2$ gives the claim. \qedhere
\end{proof}
\begin{lemma}
    A lemma right after a proof.
\end{lemma}
\begin{corollary}
    A corollary that ends in a display:
    \[ \norm{x + y} \leq \norm{x} + \norm{y}. \]
\end{corollary}
\begin{proof}
    A proof that ends in a display:
    \[ \norm{x + y}^2 \leq (\norm{x} + \norm{y})^2. \qedhere \]
\end{proof}
\begin{example}
    See Theorem~\ref{thm:cs}, \cref{thm:cs} and \eqref{eq:cs}.
\end{example}
\begin{remark}
    An unnumbered remark.
\end{remark}


\section{Floats and citations}

\Cref{tab:booktabs} and \cref{fig:tikz} float; the inequality is due to Cauchy~\cite{cauchy} and Schwarz~\cite{schwarz}.

\begin{table}[h]
    \centering
    \caption{A table with math in it.}
    \label{tab:booktabs}
    \begin{tabular}{lcr}
        \toprule
        Distribution & Mean & Variance \\
        \midrule
        $N(\mu, \sigma^2)$ & $\mu$ & $\sigma^2$ \\
        $\mathrm{Exp}(\lambda)$ & $1/\lambda$ & $1/\lambda^2$ \\
        $\mathrm{Bin}(n, p)$ & $np$ & $np(1 - p)$ \\
        \bottomrule
    \end{tabular}
\end{table}
\begin{figure}[h]
    \centering
    \begin{tikzpicture}
        \draw[->] (-0.5, 0) -- (3, 0) node[right] {$x$};
        \draw[->] (0, -0.5) -- (0, 2) node[above] {$y$};
        \draw[domain=0:2.5, smooth, thick] plot (\x, {0.25 * \x * \x});
        \node at (2, 1.6) {$y = x^2/4$};
    \end{tikzpicture}
    \caption{A TikZ picture.}
    \label{fig:tikz}
\end{figure}

\begin{thebibliography}{9}
    \bibitem{cauchy} A.-L. Cauchy, \emph{Cours d'analyse de l'\'Ecole royale polytechnique}, 1821.
    \bibitem{schwarz} H. A. Schwarz, \emph{\"Uber ein die Fl\"achen kleinsten Fl\"acheninhalts betreffendes Problem der Variationsrechnung}, 1888.
\end{thebibliography}

\end{document}
